\section{A determinant construction and its affine section}\label{sec:cone}
The first step is to place the obstruction in a tractable affine
section. We prove hyperbolicity for every admissible quadratic matrix
and identify the section exactly, including its singular boundary.

Let $Q:\R^m\to\Sym_r$ have homogeneous quadratic entries and satisfy
$Q(y)\succeq0$ for every $y$. Put $n=r+1$. For
$W=\left(\begin{smallmatrix}a&v^\top\\v&T\end{smallmatrix}\right)
\in\Sym_n$, define
\begin{equation}\label{eq:Phi}
 \Phi_y(W)=
 \begin{pmatrix}
  \tr(Q(y)T)&-v^\top Q(y)\\
  -Q(y)v&aQ(y)
 \end{pmatrix}.
\end{equation}
The polynomial used throughout the paper is
\begin{equation}\label{eq:pQ}
 p_Q(X,Z,y)=
 \det\bigl((\det X)Z-\Phi_y(\operatorname{adj}X)\bigr),
 \qquad X,Z\in\Sym_n,\quad y\in\R^m.
\end{equation}
The adjugate in this formula is the classical adjugate, so the formula
is polynomial even when $X$ is singular.

\begin{proposition}\label{prop:hyperbolic}
The polynomial $p_Q$ is homogeneous of degree $n(n+1)$ and hyperbolic
with respect to $e=(I_n,I_n,0)$, with $p_Q(e)=1$.
Fix $y$ and choose vectors $b_1,\ldots,b_n\in\R^r$ with
$Q(y)=\sum_{j=1}^n b_jb_j^\top$. Define
\[
 B_j=\begin{pmatrix}0&b_j^\top\\-b_j&0\end{pmatrix},
 \qquad
 \mathcal B=\begin{pmatrix}B_1\\\vdots\\B_n\end{pmatrix}.
\]
Then the closed hyperbolicity cone $K_Q$ of $p_Q$ is characterized,
at this value of $y$, by
\begin{equation}\label{eq:block-cone}
 (X,Z,y)\in K_Q
 \quad\Longleftrightarrow\quad
 \begin{pmatrix}I_n\otimes X&\mathcal B\\
                 \mathcal B^\top&Z\end{pmatrix}\succeq0.
\end{equation}
\end{proposition}

\begin{proof}
The entries of the matrix in \eqref{eq:pQ} have degree $n+1$, and
$p_Q(e)=1$, which proves the degree assertion. A Gram factorization
as in the statement exists since $n>r$; zero vectors may be added.
Direct multiplication gives
\[
 \Phi_y(W)=\sum_{j=1}^n B_j^\top W B_j.
\]
For every real $\eta$ we have the polynomial identity
\begin{equation}\label{eq:line-determinant}
 p_Q(X,Z,\eta y)=
 \det\begin{pmatrix}I_n\otimes X&\eta\mathcal B\\
                     \eta\mathcal B^\top&Z\end{pmatrix}.
\end{equation}
Indeed, when $X$ is invertible the right side is
\[
 (\det X)^n
 \det\bigl(Z-\eta^2\Phi_y(X^{-1})\bigr),
\]
which equals the left side by homogeneity of $Q$ and the adjugate
identity. Equality everywhere follows because both sides are
polynomials in $X,Z,\eta$.

Write $D_y(X,Z)$ for the symmetric matrix on the right of
\eqref{eq:block-cone}. Substituting $tI_n-X,tI_n-Z,-y$ in
\eqref{eq:line-determinant} gives
\[
 p_Q(te-(X,Z,y))=\det\bigl(tI_{n(n+1)}-D_y(X,Z)\bigr).
\]
Its roots are real, being the eigenvalues of a real symmetric matrix.
They are all nonnegative exactly when $D_y(X,Z)\succeq0$, proving both
hyperbolicity and \eqref{eq:block-cone}.
\end{proof}

The Gram factorization in Proposition~\ref{prop:hyperbolic} is chosen
separately for each $y$. In particular, \eqref{eq:block-cone} is not
asserted to be a linear matrix inequality jointly in $X,Z,y$.
The closed cone $K_Q$ is convex by G\aa rding's theorem
\cite{Garding1959}.

For a positive semidefinite matrix $A$, let $A^\dagger$ denote its
Moore--Penrose inverse. Define
\begin{equation}\label{eq:EQ}
 E_Q=\left\{(A,t,y):
 \begin{array}{l}
 A\in\Sym_r,\ A\succeq0,\quad \ran Q(y)\subseteq\ran A,\\[2pt]
 t\geq\tr\bigl(A^\dagger Q(y)\bigr)
 \end{array}\right\}.
\end{equation}

\begin{proposition}\label{prop:epigraph}
The set $E_Q$ is the image of the affine section
$K_Q\cap\{X=\diag(1,A):A\in\Sym_r\}$ under the projection
$(X,Z,y)\mapsto(A,Z_{00},y)$. Consequently, if $K_Q$ is a
spectrahedral shadow, then so is $E_Q$.
\end{proposition}

\begin{proof}
For a symmetric block matrix
$\left(\begin{smallmatrix}H&C\\C^\top&Z\end{smallmatrix}\right)$,
positive semidefiniteness is equivalent to
\[
 H\succeq0,\qquad \ran C\subseteq\ran H,\qquad
 Z-C^\top H^\dagger C\succeq0.
\]
To see the boundary assertion directly, positivity forces every
vector in $\ker H$ to annihilate $C$, by testing the quadratic form
on $(u,\lambda v)$. The asserted range condition then follows, and
a triangular congruence reduces the matrix to the direct sum of
$H$ and $Z-C^\top H^\dagger C$.

Apply this criterion to \eqref{eq:block-cone} with $X=\diag(1,A)$.
Its range condition says $b_j\in\ran A$ for every $j$, equivalently
$\ran Q(y)\subseteq\ran A$. Moreover,
\[
 \mathcal B^\top(I_n\otimes X)^\dagger\mathcal B
 =\Phi_y(\diag(1,A^\dagger))
 =\diag\bigl(\tr(A^\dagger Q(y)),Q(y)\bigr).
\]
It follows that every point of the affine section projects into
$E_Q$. Conversely, given $(A,t,y)\in E_Q$, take
$Z=\diag(t,Q(y))$ to obtain a point of that section. This proves
the asserted equality with no closure operation. Affine sections
and coordinate projections preserve the existence of an affine
semidefinite lift, proving the final assertion.
\end{proof}

We will construct $Q$ for which $E_Q$ has no such lift. The rest of
the argument therefore concerns \eqref{eq:EQ}, not the much larger
determinant in \eqref{eq:pQ}.
